% !TeX root = ../../../navier-stokes-manuscript.tex
% ============================================================
% 1.4 Silver Standard Proof
% ============================================================
\subsection{Silver Standard Proof}\label{sub:SilverStandardProof}

Finding an estimate strong enough to finish the Gold-standard proof is genuinely difficult.
No one has closed that estimate for arbitrary smooth three-dimensional data, and once you dig your teeth into the problem, it becomes tempting to look for another way of proving smoothness.

Assume, for a moment, that the problem has already been solved by the Gold-standard method.

That would mean every fluid produced by the original Navier--Stokes problem remains smooth.
Its velocity, pressure, viscosity, incompressibility, and nonlinear transport would all remain in force throughout the evolution.
Now imagine a blow-up scenario, or any other outcome that breaks smoothness.
Under the Gold assumption, that outcome cannot be reached while every requirement of the original problem remains intact.
At least one requirement must have been lost, ignored, or replaced along the way.
The resulting nonsmooth object would no longer describe the same fluid that began from \(u_0\).

So in short, the idea becomes:

If the original Navier--Stokes equations always produce a smooth fluid, then a nonsmooth outcome cannot be a result of those same equations.

In symbols,
\[
  \text{Navier--Stokes}
  \Longrightarrow
  \text{smooth}
\]
has the contrapositive
\[
  \text{not smooth}
  \Longrightarrow
  \text{not Navier--Stokes}.
\]

The Gold assumption revealed the second implication.
A Silver proof has to establish it directly by showing that every alleged nonsmooth outcome loses at least one requirement of the original fluid.
That is why the Silver Standard is a proof by contrapositive.

The beauty of the Silver method is the question it lets us ask: is this nonsmooth object still Navier--Stokes?
Gold searches for the estimate that keeps the solution smooth through every finite time.
Silver tests every nonsmooth proposal against the properties that define the original fluid.
Its advantage is a simpler question.
Its difficulty is the number of possible nonsmooth outcomes that question appears to create.
A single principle that reaches every case would make the Silver method usable.
